\documentclass[12pt,a4paper]{report}

% --- Gói ngôn ngữ và mã hóa tiếng Việt ---
\usepackage[utf8]{inputenc}
\usepackage[vietnamese]{babel}
\usepackage[T5]{fontenc}

% --- Canh lề và định dạng trang ---
\usepackage[top=2.5cm, bottom=2.5cm, left=3cm, right=2cm]{geometry}
\usepackage{fancyhdr}
\setlength{\headheight}{15pt}
\usepackage{titlesec}

% --- Bảng, hình ảnh, toán học và ký hiệu ---
\usepackage{amsmath, amssymb}
\usepackage{graphicx}
\usepackage{booktabs}
\usepackage{tabularx}
\usepackage{float}
\usepackage{caption}
\usepackage{subcaption}

% --- Định dạng code mã nguồn ---
\usepackage{listings}
\usepackage{xcolor}

\definecolor{codegreen}{rgb}{0.13,0.55,0.13}
\definecolor{codegray}{rgb}{0.5,0.5,0.5}
\definecolor{codepurple}{rgb}{0.58,0,0.82}
\definecolor{backcolour}{rgb}{0.96,0.97,0.98}

\lstdefinestyle{mystyle}{
    backgroundcolor=\backcolour,   
    commentstyle=\color{codegreen},
    keywordstyle=\color{blue}\bfseries,
    numberstyle=\tiny\color{codegray},
    stringstyle=\color{codepurple},
    basicstyle=\ttfamily\footnotesize,
    breakatwhitespace=false,         
    breaklines=true,                 
    captionpos=b,                    
    keepspaces=true,                 
    numbers=left,                    
    numbersep=5pt,                  
    showspaces=false,                
    showstringspaces=false,
    showtabs=false,                  
    tabsize=2
}
\lstset{style=mystyle}

% --- Liên kết và màu sắc tài liệu ---
\usepackage{hyperref}
\hypersetup{
    colorlinks=true,
    linkcolor=blue!70!black,
    filecolor=magenta,      
    urlcolor=cyan!70!black,
    citecolor=green!50!black,
    pdftitle={Báo cáo dự án IoT: FloodGuard - Hệ thống chống ngập tự động},
    pdfauthor={Nhóm tác giả},
}

% --- Header và Footer ---
\pagestyle{fancy}
\fancyhf{}
\fancyhead[L]{\small\slshape FloodGuard - Hệ Thống Chống Ngập Tự Động}
\fancyhead[R]{\small\slshape ESP32 IoT Project}
\fancyfoot[C]{\thepage}
\renewcommand{\headrulewidth}{0.4pt}
\renewcommand{\footrulewidth}{0.4pt}

% ==================== BẮT ĐẦU TÀI LIỆU ====================
\begin{document}

% --- TRANG BÌA ---
\begin{titlepage}
    \centering
    \vspace*{0.3cm}
    
    {\fontsize{15pt}{19pt}\selectfont\textbf{ĐẠI HỌC QUỐC GIA HÀ NỘI}\par}
    \vspace{0.15cm}
    {\fontsize{16pt}{20pt}\selectfont\textbf{TRƯỜNG ĐẠI HỌC CÔNG NGHỆ}\par}
    \vspace{0.4cm}
    \hrule height 1.5pt
    \vspace{0.8cm}
    
    \includegraphics[width=5.2cm]{figures/logo-uet.jpg}\par
    \vspace{0.9cm}
    
    {\Large\textbf{BÁO CÁO BÀI TẬP LỚN}\par}
    \vspace{0.2cm}
    {\large\textbf{MÔN: LẬP TRÌNH NHÚNG VÀ THỜI GIAN THỰC}\par}
    \vspace{0.8cm}
    
    {\textbf{ĐỀ TÀI:}\par}
    \vspace{0.2cm}
    {\fontsize{13pt}{17pt}\selectfont\bfseries FLOODGUARD: HỆ THỐNG GIÁM SÁT VÀ CHỐNG NGẬP TỰ ĐỘNG THÔNG MINH\par}
    \vspace{0.3cm}
    {\small\slshape Ứng dụng Vi điều khiển ESP32, Thuật toán FSM \& Web Dashboard Thời Gian Thực\par}
    
    \vspace{1.6cm}
    
    \begin{minipage}{0.92\textwidth}
        \normalsize
        \renewcommand{\arraystretch}{1.35}
        \begin{tabular}{@{}ll@{}}
            \textbf{Học phần:}            & Lập trình nhúng và thời gian thực \\
            \textbf{Mã học phần:}         & 2627I\_INT3108\_1 \\
            \textbf{Sinh viên thực hiện:} & Nguyễn Văn B - MSSV: 202xxxxx \\
            \textbf{Học kỳ / Năm học:}    & Học kỳ 1 -- 2026-2027 \\
        \end{tabular}
    \end{minipage}
    
    \vfill
    {\large Hà Nội, Tháng 10 Năm 2026}\par
    \vspace*{0.3cm}
\end{titlepage}

% --- MỤC LỤC ---
\pagenumbering{roman}
\tableofcontents
\newpage
\listoffigures
\listoftables
\newpage

% --- NỘI DUNG CHÍNH ---
\pagenumbering{arabic}

\chapter{Giới Thiệu Đề Tài}

\section{Đặt Vấn Đề}
Hiện tượng ngập lụt đô thị và sạt lở, tràn đê vào mùa mưa bão đang là một trong những thách thức nghiêm trọng nhất tại các đô thị lớn ở Việt Nam cũng như trên thế giới. Tình trạng này gây thiệt hại nặng nề về kinh tế, tài sản, cản trở giao thông và tiềm ẩn nhiều rủi ro về sức khỏe, tính mạng con người.

Các giải pháp truyền thống như điều tiết thủ công, đo đạc bằng mắt thường hoặc vận hành trạm bơm tại chỗ bộc lộ nhiều nhược điểm:
\begin{itemize}
    \item \textbf{Độ trễ phản ứng cao:} Khi nước dâng nhanh đột ngột vào ban đêm, con người không kịp có mặt để vận hành máy bơm hay mở cống xả.
    \item \textbf{Thiếu thông tin cảnh báo sớm:} Người dân trong khu vực không được cập nhật kịp thời về tốc độ và mức độ nước dâng.
    \item \textbf{Nguy cơ hỏng hóc thiết bị:} Việc bật/tắt thiết bị công suất lớn liên tục do mực nước dập dềnh gây giảm tuổi thọ của bơm và van xả.
\end{itemize}

\section{Mục Tiêu Dự Án FloodGuard}
Dự án \textbf{FloodGuard} hướng tới việc xây dựng một hệ thống IoT giám sát và chống ngập tự động, nhỏ gọn, hoạt động bền bỉ và phản ứng tức thời. Hệ thống tập trung vào các mục tiêu then chốt:
\begin{enumerate}
    \item Tự động đo lường mực nước và lượng mưa liên tục theo thời gian thực với độ chính xác cao.
    \item Tự động ra quyết định kích hoạt máy bơm nước và góc mở van thoát lũ theo thuật toán máy trạng thái hữu hạn (\textit{Finite State Machine - FSM}) kết hợp độ trễ trễ (\textit{Hysteresis}).
    \item Phát tín hiệu cảnh báo đa tầng thông qua hệ thống đèn LED trạng thái và còi hú báo động âm lượng lớn.
    \item Cung cấp giao diện Web Dashboard trực quan, cho phép người quản lý theo dõi và can thiệp điều khiển thủ công (\textit{Manual Override}) từ xa qua mạng Wi-Fi.
\end{enumerate}


\chapter{Kiến Trúc Phần Cứng \& Phân Công Chân}

\section{Tổng Quan Kiến Trúc Hệ Thống}
Hệ thống FloodGuard gồm 4 khối chức năng chính:
\begin{enumerate}
    \item \textbf{Khối xử lý trung tâm:} Vi điều khiển ESP32 DevKit V1 với chip lõi kép Xtensa 32-bit, tích hợp kết nối Wi-Fi 802.11 b/g/n, đóng vai trò trạm thu thập tín hiệu và máy chủ Web Server thời gian thực.
    \item \textbf{Khối cảm biến (Sensors):} Cảm biến siêu âm HC-SR04 đo mực nước hồ thu và chiết áp/cảm biến mưa đo độ ẩm, lượng nước mưa.
    \item \textbf{Khối chấp hành (Actuators):} Động cơ Servo SG90 điều khiển góc mở van xả lũ (0° – 90° – 180°) và Module Relay 5V kích đóng/ngắt máy bơm thoát nước.
    \item \textbf{Khối cảnh báo tại chỗ (Signaling):} Cụm 3 đèn LED màu (Xanh, Vàng, Đỏ) và còi báo động Active Buzzer 5V.
\end{enumerate}

\section{Giải Pháp Điều Khiển Qua Web Thay Thế Nút Nhấn Vật Lý}
Trong phiên bản thiết kế cải tiến, hệ thống \textbf{loại bỏ hoàn toàn các nút bấm cơ học vật lý} và chuyển sang các nút điều khiển ảo trực tiếp trên giao diện Web Dashboard. Giải pháp này mang lại nhiều ưu điểm kỹ thuật vượt trội:
\begin{itemize}
    \item \textbf{Độ bền cao trong môi trường ẩm ướt:} Môi trường chống ngập bão lũ có độ ẩm rất cao, nút nhấn vật lý rất dễ bị oxy hóa tiếp điểm hoặc rò rỉ nước gây chập mạch.
    \item \textbf{An toàn cho người vận hành:} Kỹ thuật viên không cần tiếp cận trực tiếp khu vực hồ thu hay trạm bơm đang ngập sâu, mà có thể điều khiển từ xa thông qua điện thoại hoặc máy tính kết nối Wi-Fi.
    \item \textbf{Tiết kiệm chân GPIO vi điều khiển:} Tiết kiệm được các chân GPIO ngắt ngoài để phục vụ mở rộng thêm cảm biến trong tương lai.
\end{itemize}

\section{Bảng Phân Công Chân GPIO}
\begin{table}[H]
\centering
\caption{Bảng phân công chân phần cứng trên ESP32}
\label{tab:gpio_mapping}
\begin{tabularx}{\textwidth}{l l l X}
\toprule
\textbf{Thiết bị / Module} & \textbf{Chân tín hiệu} & \textbf{GPIO ESP32} & \textbf{Chức năng} \\
\midrule
LED Xanh & Anode (+) & \texttt{GPIO 25} & Cảnh báo mức an toàn (NORMAL) \\
LED Vàng & Anode (+) & \texttt{GPIO 26} & Cảnh báo mức dâng sớm (WARNING) \\
LED Đỏ & Anode (+) & \texttt{GPIO 27} & Cảnh báo mức nguy cấp (DANGER / CRITICAL) \\
Buzzer 5V & (+) & \texttt{GPIO 19} & Còi hú báo động khẩn cấp \\
HC-SR04 (Trig) & Trigger & \texttt{GPIO 18} & Phát xung siêu âm kích hoạt \\
HC-SR04 (Echo) & Echo & \texttt{GPIO 5} & Nhận xung phản hồi phản xạ từ mặt nước \\
Cảm biến mưa & Analog Out & \texttt{GPIO 34} & Đo cường độ mưa / độ ẩm \\
Servo SG90 & PWM & \texttt{GPIO 13} & Điều khiển độ mở van thoát nước \\
Relay 1 kênh & IN & \texttt{GPIO 23} & Kích hoạt đóng/ngắt máy bơm mini \\
\bottomrule
\end{tabularx}
\end{table}


\chapter{Thuật Toán Điều Khiển \& Vận Hành}

\section{Máy Trạng Thái Hữu Hạn (FSM)}
Để đảm bảo hệ thống phản ứng tin cậy và không bị nhiễu động, hệ thống sử dụng mô hình Máy trạng thái hữu hạn (\textit{Finite State Machine - FSM}) với 4 trạng thái cốt lõi:

\begin{itemize}
    \item \textbf{Trạng thái 1: [NORMAL] (Mực nước 0\% -- 20\%):} 
    Mực nước ở ngưỡng an toàn tuyệt đối. Đèn LED xanh sáng liên tục, máy bơm ở trạng thái OFF, van xả đóng hoàn toàn (góc 0°), còi báo động tắt.
    
    \item \textbf{Trạng thái 2: [WARNING] (Mực nước 21\% -- 40\%):} 
    Nước bắt đầu dâng vào hồ thu. Đèn LED vàng sáng cảnh báo sớm cho kỹ thuật viên và người dân, các cơ cấu bơm và xả vẫn ở chế độ chờ.
    
    \item \textbf{Trạng thái 3: [DANGER] (Mực nước 41\% -- 70\%):} 
    Mực nước vượt ngưỡng kiểm soát thông thường. Đèn LED đỏ bật sáng, Relay đóng kích hoạt máy bơm nước mini chạy hết công suất, động cơ servo xoay 90° để mở 50\% khẩu độ van xả lũ.
    
    \item \textbf{Trạng thái 4: [CRITICAL] (Mực nước 71\% -- 100\%):} 
    Tình trạng ngập lụt nghiêm trọng, đe dọa tràn hồ. Đèn LED đỏ nhấp nháy liên tục với tần số cao, còi hú dồn dập, máy bơm tiếp tục hoạt động, động cơ servo xoay góc tối đa 180° để mở 100\% van thoát nước.
\end{itemize}

\section{Nguyên Lý Chống Dao Động (Hysteresis Band)}
Trong thực tế, mặt nước trong hồ thu thường có hiện tượng gợn sóng hoặc dao động liên tục xung quanh các mốc ngưỡng chuyển trạng thái (ví dụ mức 40\% hoặc 70\%). Nếu không có cơ chế bù trừ, máy bơm và van xả sẽ bị bật/tắt liên tục trong thời gian ngắn, gây cháy cuộn hút relay hoặc hỏng động cơ.

FloodGuard áp dụng dải trễ (\textit{Hysteresis}):
\begin{itemize}
    \item \textbf{Khi nước dâng:} Hệ thống chỉ kích hoạt Bơm và Van khi nước đạt hoặc vượt quá ngưỡng 41\%.
    \item \textbf{Khi nước rút:} Bơm và Van vẫn duy trì trạng thái xả nước cho tới khi mực nước hạ hẳn xuống dưới \textbf{40\%} mới bắt đầu ngắt bơm đóng van, và phải hạ sâu dưới \textbf{20\%} thì hệ thống mới khôi phục lại trạng thái NORMAL.
\end{itemize}

\chapter{Giao Diện Web Dashboard Giám Sát \& Can Thiệp Từ Xa}

\section{Kiến Trúc Web Server Trên ESP32}
ESP32 đóng vai trò là một HTTP Web Server hoạt động trên cổng tiêu chuẩn 80. Trang Dashboard được xây dựng theo phong cách hiện đại Dark Glassmorphism, tối ưu hóa giao diện phản hồi (\textit{Responsive Design}) cho cả trình duyệt điện thoại và máy tính để bàn.

\section{Các Thành Phần Giám Sát Thời Gian Thực}
\begin{itemize}
    \item \textbf{Thanh đo mực nước (Water Level Gauge):} Hiển thị phần trăm dung tích chứa (0\% -- 100\%) và khoảng cách cm thực tế từ cảm biến siêu âm tới mặt nước, kèm thanh màu cảnh báo trực quan.
    \item \textbf{Cảm biến lượng mưa:} Cập nhật trạng thái trời mưa (\textit{CÓ MƯA} / \textit{KHÔNG MƯA}) và giá trị đo thô ADC từ cảm biến.
    \item \textbf{Thẻ trạng thái hệ thống:} Huy hiệu hiển thị trạng thái FSM hiện tại (\textit{NORMAL}, \textit{WARNING}, \textit{DANGER}, \textit{CRITICAL}) và chế độ hoạt động (\textit{AUTO} hoặc \textit{MANUAL}).
    \item \textbf{Giám sát thiết bị chấp hành:} Hiển thị trạng thái hoạt động thực tế của Máy bơm (ĐANG BẬT / ĐÃ TẮT), góc mở của Van xả (0\% -- 50\% -- 100\%) và trạng thái Còi hú.
\end{itemize}

\section{Hệ Thống Nút Bấm Ảo \& Điều Khiển Can Thiệp Từ Xa}
Thay vì sử dụng các nút bấm vật lý trên bo mạch (dễ hỏng trong môi trường nước và mưa bão), giao diện Web Dashboard tích hợp cụm nút bấm điều khiển số hóa trực tiếp:

\begin{itemize}
    \item \textbf{Nút Đổi Chế Độ (\texttt{mode\_toggle}):} Cho phép chuyển đổi linh hoạt giữa chế độ vận hành tự động (\textit{AUTO}) và điều khiển thủ công (\textit{MANUAL}).
    \item \textbf{Nút Bật/Tắt Máy Bơm (\texttt{pump\_toggle}):} Trong chế độ MANUAL, kỹ thuật viên có thể chủ động kích hoạt hoặc ngắt máy bơm mini ngay tức thì.
    \item \textbf{Cụm nút điều chỉnh Van xả (\texttt{valve\_0}, \texttt{valve\_50}, \texttt{valve\_100}):} Điều khiển van xả theo 3 mức khẩu độ linh hoạt: Đóng hoàn toàn (0\%), Mở 50\% (góc 90°) hoặc Mở thoát lũ tối đa 100\% (góc 180°).
\end{itemize}

\section{Giao Thức Truyền Nhận Dữ Liệu HTTP REST API}
Giao diện Web liên tục đồng bộ dữ liệu với ESP32 thông qua kỹ thuật bất đồng bộ AJAX/Fetch API mỗi giây mà không cần tải lại toàn bộ trang web:
\begin{itemize}
    \item \textbf{\texttt{GET /api/status}:} Trả về chuỗi định dạng JSON chứa toàn bộ dữ liệu đo đạc (khoảng cách, lượng mưa, trạng thái máy bơm, góc van, chế độ hoạt động).
    \item \textbf{\texttt{GET /api/control?action=\{cmd\}}:} Nhận lệnh điều khiển từ các nút bấm ảo trên giao diện để cập nhật biến trạng thái bên trong vi điều khiển ESP32.
\end{itemize}


\chapter{Kết Luận \& Hướng Phát Triển}

\section{Kết Quả Đạt Được}
Dự án \textbf{FloodGuard} đã hoàn thành đầy đủ các mục tiêu thiết kế ban đầu:
\begin{enumerate}
    \item Thiết kế và mô phỏng hoàn chỉnh hệ thống trên nền tảng Wokwi kết hợp vi điều khiển ESP32.
    \item Thuật toán FSM và dải trễ Hysteresis vận hành ổn định, loại bỏ triệt để hiện tượng rung giật đóng cắt thiết bị.
    \item Giao diện Web Dashboard trực quan, truyền nhận dữ liệu nhanh chóng qua Wi-Fi thời gian thực.
\end{enumerate}

\section{Hướng Phát Triển Trong Tương Lai}
\begin{itemize}
    \item Kết nối truyền thông qua giao thức MQTT và dịch vụ điện toán đám mây (AWS IoT Core / ThingsBoard).
    \item Tích hợp thuật toán học máy (Machine Learning) nhẹ để dự báo đỉnh triều cường và lưu lượng mưa trước 1--2 giờ.
    \item Trang bị thêm nguồn pin dự phòng (UPS/Solar) đảm bảo hệ thống không bị gián đoạn khi mất điện lưới do bão lũ.
\end{itemize}



\end{document}
